\documentclass[12pt,a4paper,oneside]{article}
\usepackage[brazil]{babel}
\usepackage[utf8]{inputenc}
\usepackage{olymp}

\setcounter{problem}{4}

\begin{document}

\begin{problem}{Distância}{distancia}{2}
 
\textit{Googlemaps}, \textit{ViaMichelin} e outros sites do gênero calculam rotas entre duas localizações. \textit{ViaMichelin} oferece opções de rota mais curta (em distância), mais rápida (em tempo), mais barata (incluindo pedágios) e em alguns países também a opção de rota mais ``bonita'', além de informar consumo de combustível.


Você deve fazer um programa que calcula a menor distância entre duas localizações, mas de forma mais simples: as localizações serão pontos no plano cartesiano, e a menor distância é a distância entre eles em linha reta. A figura abaixo mostra, à esquerda, um exemplo com dois pontos, $P_1$ e $P_2$, de coordenadas $(2,1)$ e $(6,4)$, respectivamente. A da direita mostra a distância entre eles, que vale $5$.

\begin{center}
\includegraphics[scale=0.25]{images/exercise_100_0}
\end{center}

%\Notes
%
%Observações, se tiver.

\InputFile

A entrada terá apenas um caso de teste.

O caso de teste é composto por quatro valores reais, $X_1, Y_1, X_2, Y_2$, sendo $(X_1,Y_1)$ as coordenadas cartesianas do primeiro ponto e $(X_2,Y_2)$ as do segundo ponto.. Restrições: $0 < X_1,Y_1,X_2,Y_2 \leq 1000$.


\OutputFile

Seu programa deve gerar apenas uma linha de saída, contendo um valor real, formatado com duas casas decimais, indicando a distância entre os dois pontos.

Não se esqueça de quebrar a linha após a impressão do valor.

\Example

\begin{example}
\exmp{
2 1 6 4
}{
5.00
}%
\end{example}

\begin{example}
\exmp{
10 20 70 100
}{
100.00
}%
\end{example}

\begin{example}
\exmp{
1.1 30 13.6 25
}{
13.46
}%
\end{example}


%Comentários sobre o exemplo, se necessário.

\end{problem}

\end{document}
